\section{26 сентября 1914 г.}

\lp{20260604_193445066}

Дорогой мой! У меня такое полное доверие к вам и такая невозможность быть не с
вами, что и речи быть не может, чтобы мы разошлись в коренном. Тем более, что
ведь и мое отношение к н. \rem{немцам} не такое, чтобы какие-нибудь факты могли
во мне все перевернуть. Но дело в том, что это надо теперь понять, согласовать
и вообще проделать очень трудную работу, чтобы продолжать дальше жить. Коля от
этой войны так невыразимо страдает, что я одно время прямо-таки опасалась за
его здоровье. И Коле и мне теперь необходимо было бы поговорить об этом с вами.
Как объяснить все-таки некоторые поступки н.? Поведение с Бельги\rem{...}?
Зачем расстреливали и женщин и детей? Это не из русских газет, а из писем
известных бельгийцев. Вообще много

\lp{20260604_193503074}

\noindent
надо объяснить и понять, чтобы принять. Я не могу отказаться от своей любви к
н. и потому ужасно страдаю от некоторых вопросов. Вопрос -- верна ли себе
действительно современная Германия? Одобрили ли бы все Гёте? Ужасно, что
невозможно все писать. Поверьте мне только, мой дорогой, это и в н.
\rem{немецких} газетах не все печатается. В этом отношении великолепны
рассказы(печатаются в «Русском Вече») Станиславского, в которых он с полной
справедливостью описывает все, что с ними было, и как выступили и положительные
и отрицательные стороны. Мне даже больно писать об этом, да и нельзя. Смысла
тоже нет, а гораздо важнее было бы поговорить. Когда то это теперь
будет\textellipsis

Но я все время радуюсь, что вы в стороне от тревог. Все поручения исполню в
точности. Письмо Смирнову отнесу сама. Фридрих пошлю по почте. (Между прочим,
неужели Ядвига решила вообще никак не относиться ко мне? Меня этот вопрос
все-таки интересует. А впрочем, если вам это все равно, то тогда и Бог с ним!).

\lp{20260604_193456434}

Я рада, что Мариетта вполне согласна с вами. Мне бы с ней увидаться поскорее!
